% ============================================
% Assistant Professor Cover Letter – Hampden–Sydney College
% ============================================

\documentclass[11pt,letterpaper]{letter}

\usepackage{geometry}
\usepackage{microtype}
\usepackage[hidelinks]{hyperref}

\geometry{left=25mm,right=25mm,top=25mm,bottom=25mm}

\signature{Decory Edwards}
\address{
Department of Economics \\
Johns Hopkins University \\
Baltimore, MD 21218 \\
\texttt{dedwar65@jh.edu}
}

\begin{document}

\begin{letter}{Search Committee \\
Department of Economics and Business \\
Hampden--Sydney College \\
Hampden--Sydney, VA \\ USA}

\opening{Dear Members of the Search Committee,}

I am writing to apply for the tenure-track Assistant Professor position in Economics and Business at Hampden--Sydney College. I am a Ph.D.\ candidate in Economics at Johns Hopkins University, advised by Professors Christopher D.~Carroll, Nicholas W.~Papageorge, and Stelios Fourakis, and I expect to receive my degree in May~2026. My research fields are macroeconomics, wealth and income dynamics, and computational economics.

My research agenda focuses on understanding the determinants of wealth inequality and the mechanisms through which heterogeneity in returns to wealth shapes aggregate outcomes. In my job-market paper, I estimate the degree of heterogeneity in individual portfolio returns using micro-data from the Survey of Consumer Finances and simulate a structural model in which households face idiosyncratic return risk. I show that allowing for persistent heterogeneity in returns substantially improves the model’s ability to match the empirical distribution of wealth, offering a tractable way to connect micro-level return dispersion to macroeconomic inequality.

In a second project, I use the Health and Retirement Study to examine the relationship between interpersonal trust and portfolio performance. Building on the literature that finds a hump-shaped relationship between trust and income, I show that a similar relationship holds for individual returns to wealth measured in the HRS data, highlighting a behavioral channel through which social attitudes shape saving and investment outcomes. This work also provides new estimates of the persistent component of returns to net worth, which motivate the calibration of heterogeneity in my structural model.

\textbf{Teaching.} I am strongly committed to undergraduate teaching and to the teacher–scholar model emphasized at Hampden--Sydney College. I have experience teaching and supporting courses in \textit{Principles of Macroeconomics}, \textit{Intermediate Macroeconomic Theory}, and quantitative methods. My teaching emphasizes economic intuition, careful reasoning, and the disciplined use of data to evaluate policy and business decisions. At Hampden--Sydney, I would be well prepared to teach courses such as intermediate microeconomics, econometrics and forecasting, applied public policy, and research methods, while mentoring students through independent projects and senior research. I value close interaction with students and see undergraduate research as a natural extension of classroom learning.

Hampden--Sydney’s mission as a residential liberal arts college dedicated to forming thoughtful, responsible citizens is especially appealing to me. I am drawn to the College’s emphasis on rigorous instruction, close faculty–student engagement, and the integration of quantitative reasoning with ethical and civic considerations. I would welcome the opportunity to contribute to a collegial department, to participate actively in service, and to support the broader educational mission of the College.

Thank you very much for your time and consideration. I would welcome the opportunity to discuss my application further.

\closing{Sincerely,}

\end{letter}
\end{document}
